\documentclass[a4paper, 12pt]{article}
\usepackage[utf8]{inputenc}
\usepackage[margin=1in]{geometry}
\usepackage{setspace}
\usepackage{xcolor}
\usepackage{hyperref}
\usepackage{lipsum}
\usepackage{fancyhdr}
\usepackage[para]{footmisc}

\renewcommand{\headrulewidth}{0pt}
\fancyfoot[R]{\thepage}
\singlespacing
\usepackage{titlesec}
\titleformat{\section}
  {\normalfont\large\bfseries}{\thesection}{1em}{}
\titleformat{\subsection}
  {\normalfont\normalsize\bfseries}{\thesubsection}{1em}{}
\titleformat{\subsubsection}
  {\normalfont\small\bfseries}{\thesubsubsection}{1em}{}
\titlespacing*{\section}{0pt}{1.5ex plus 1ex minus .2ex}{1ex plus .2ex}
\titlespacing*{\subsection}{0pt}{1.5ex plus 1ex minus .2ex}{0.8ex plus .2ex}
\titlespacing*{\subsubsection}{0pt}{1ex plus 1ex minus .2ex}{0.5ex plus .2ex}
\renewcommand{\thesubsection}{\arabic{subsection}}

\pagestyle{fancy}
\fancyhf{}
\fancyhead[L]{IK2221 VT26}
\fancyhead[R]{Lorenzo Deflorian (ldef@kth.se)}
%\fancyfoot[R]{\thepage}

\begin{document}

\section*{Metron: NFV Service Chains at the True Speed of the Underlying Hardware}
\subsection{Paper summary}
% Paper summary (for E level)
%% Problem being tackled and why it's hard
%% Key idea behind the approach
%% A bit more detail about the approach
%% The way the idea is evaluated, plus the key findings

This paper presents Metron, which is an NFV platform designed to improve the performance of service chains. The main issue they try to solve is that in typical NFV systems, packets are often moved between CPU cores multiple times during processing, which creates overhead and reduces performance. This is especially bad for latency and throughput.

Metron tries to fix this by combining network hardware and servers in a smarter way. The idea is to classify packets early using tags, and then use the network hardware to send each packet directly to the correct CPU core. This way, packets don't need to be moved between cores anymore, at least for the stateful parts of the processing, which is where most of the cost usually is.

The system is managed by a controller that analyzes the service chain and divides traffic into classes. It then decides what part of the processing can be done in the network (stateless) and what needs to stay on the server (stateful). The controller also handles load balancing by splitting or merging traffic classes when needed, even if this process is not always super clean.

They evaluate Metron and show that it can handle very high throughput (up to 40 Gbps for deep inspection and close to 100 GbE speeds in some cases). It also performs significantly better than previous systems like OpenBox, with improvements in both efficiency and latency, even if the comparison is a bit specific to their setup.

\subsection{Contributions}
% Top 3 contributions of the paper  (extra, for C level)
%% Tie the contributions with impact.
%% The first paper to propose… | Makes it possible to… 
%% Presents a new technique for… | Dramatically improves performance of … 
%% Offers largely improved reliability… | Saves a large fraction of power for…

The first contribution is the idea of combining programmable network devices with commodity servers to improve NFV performance. By doing part of the processing in the network, Metron avoids unnecessary inter-core communication, which is one of the main bottlenecks in these systems.

The second contribution is the tagging-based dispatching mechanism. Packets are classified early and then sent directly to the right core, instead of being repeatedly redirected in software. This reduces overhead and makes the system more efficient, while still keeping some level of flexibility in how traffic is handled.

The third contribution is the controller and resource management strategy. The system can adapt to changes in traffic by splitting or merging traffic classes and updating the rules accordingly. This makes it more practical for real deployments, even if the decisions are based on heuristics and not always optimal.

\subsection{Problems}
% Most glaring problem(s) with the work (up to 3) (extra, for A level)
%% Does the approach scale as the number of users increases?
%% Scaling: non-linear increase in resources | How does it work under a failure?
%% Does it miss some important cases? | Does it have too strong assumptions?
%% Does it have unrealistic deployment requirements?

One concern is that Metron relies quite a lot on programmable hardware and coordination between different components (switches, NICs, servers). Even if the paper says it can work with commodity hardware, in practice you still need specific features like tag-based dispatching, which might not always be available. So deployment could be limited in some environments.

Another issue is how state is handled when traffic classes change. When classes are split or merged, the system may duplicate or move state, but some redundant state can remain for a while. The paper mentions this, but doesn't really solve it completely, so it feels more like a practical workaround than a clean solution.

A third point is scalability of the control logic. The paper itself says that tracking very fine-grained traffic classes would not be practical, so it groups them and uses heuristics to decide when to split or merge. This works in their experiments, but it might be less accurate in very large or unpredictable scenarios, where traffic patterns change a lot.

\end{document}